%!TEX root= ../../note
\subsection{Quantifiers and Quantified Statements}
Quantifiers specify how many elements of a domain must satisfy a condition.
When quantifiers are nested, their order determines which choices may
depend on earlier ones.
\subsubsection{Universal Quantifier}
$\forall$ means ``for all'', and it is called the universal quantifier.
The statement $\forall x \in S, P(x)$ is true when $P(x)$ is true for \emph{every} element $x$ of $S$.

\begin{example}
    ``Every integer $n \geq 5$ satisfies $2^n > n^2$'' can be written as
    \begin{equation*}
        \forall n \in \set{n \in \Z : n \geq 5},\ 2^n > n^2.
    \end{equation*}
    Checking $n = 5, 6, 7$ is not a proof: a universal statement must hold for every $n$ in the domain.
    (The domain matters: the statement fails at $n = 3$, since $8 < 9$.)
\end{example}

The negation of a universally quantified statement is
\begin{align*}
    \neg \paren{\forall x \in S, P(x)} \equiv \paren{\exists x \in S, \neg P(x)}.
\end{align*}
So a single counterexample is enough to show that a universal statement is false.

\subsubsection{Existential Quantifier}
$\exists$ means ``there exists'', and it is called the existential quantifier.
The statement $\exists x \in S, P(x)$ is true when $P(x)$ is true for \emph{at least one} $x \in S$.
\begin{example}
    \begin{align*}
        \exists a \in \Z, a^2 + 29a + 209 \leq 0.
    \end{align*}
    This is true: one witness is enough, and $a = -14$ gives $196 - 406 + 209 = -1 \leq 0$.
\end{example}

The negation of an existentially quantified statement (such as $\exists x \in S, P(x)$) is
\begin{equation*}
    \neg \paren{\exists x \in S, P(x)} \equiv \paren{\forall x \in S, \neg P(x)}.
\end{equation*}

\begin{remark}
    A quantified statement contains four parts:
    \begin{itemize}
        \item variable;
        \item quantifier;
        \item domain;
        \item an open sentence involving the variable (that is either true or false whenever a value of the variable is chosen from the domain).
    \end{itemize}
\end{remark}

\subsubsection{Nested Quantifiers}
Consider the statement
\begin{equation*}
    \forall x \in \R, \exists y \in \R, x > y.
\end{equation*}

This statement is true: given any $x \in \R$, we may choose $y = x - 1$, and then $x > y$.
Note that $y$ is allowed to depend on $x$.

\begin{remark}
    Reversing the quantifiers changes the statement:
    $\exists y \in \R, \forall x \in \R, x > y$.
    It says there is a single real number $y$ which is smaller than every real number,
    including $y$ itself. This is false: taking $x = y$ gives $y > y$.
\end{remark}

\medskip\noindent
To negate a nested quantified statement, negate one quantifier at a time from the outside in:
each $\forall$ becomes $\exists$, each $\exists$ becomes $\forall$, and the open sentence is negated.
For example,
\begin{align*}
    &\neg \paren{\forall x \in S, \exists y \in T, P(x, y)}\\
    &\qquad\equiv \exists x \in S, \forall y \in T, \neg P(x, y).
\end{align*}

\subsubsection*{In practice}
\begin{itemize}
    \item \textbf{Reading quantifiers.} State the domain and read from left
    to right. In $\forall x\in S,\ \exists y\in T,\ P(x,y)$, the choice of
    $y$ may depend on $x$; reversing the quantifiers requires one $y$ that
    works for every $x$.
    \item \textbf{Negating a statement.} Keep the domains and quantifier
    order, swap each $\forall$ with $\exists$, and negate the final condition.
    One counterexample refutes a universal statement; refuting an existential
    statement requires ruling out every possible witness.
\end{itemize}
